\paragraph{\trRelative\trBraneStructure。}\label{sub:relative-brane-structures}
设 $\pi\colon E\to S$ 是一个\trLefschetzFibration（我们将先讨论 $S$ \trCompact{}的情形，而把带\trStriplikeEnd{}的版本留到后面再做）。设 $E$ 带有相对二次复\trVolume{}形式 $\eta^2_{E/S}$，并令 $\alpha_{E/S}$ 为\trRelative\trSquaredPhase{}映射。再设 $F$ 是 $E$ 上的一个\trLagrange\trBoundaryCondition。$F$ 的一个\textbf{\trRelative\trGrading}（relative grading）是一个\trLift{}$\alpha^{\#}\colon F\to\RR$，它是 $x\mapsto\alpha_{E/S}(TF_x^v)$ 的\trLift。$F$ 上的一个\textbf{\trRelative\trBraneStructure}（relative brane structure）是一个对 $(\alpha^{\#},P^{\#})$，由一个\trRelative\trGrading{}与一个\textbf{\trRelative\trPinStructure}（relative Pin structure）组成，后者不过是\trVectorBundle{}$TF^v\to F$ 上的一个\trPinStructure{}$P^{\#}$ 而已。这类结构的存在性\trObstruction{}可以如下看出：在每条\trBoundaryCircle{}$C$ 上取一点 $z$。那么，每个 $F_{z_j}\subset E_z$ 都必须容许一个相对于 $\eta^2_{E_z}$ 的\trGrading\footnote{译注：原书此处印作 $F_{z_j}$；疑应为 $F_z$（上文取点记作 $z$，本段下文亦用 $F_z$）；此处照原书排印，存疑备案。——译者注}；此外，若 $h_C$ 是沿边界一周的\trMonodromy——按\cref{sub:relative-volume-forms}中的方式配备其自然\trGrading——那么作为\trGraded\trLagrange\trSubmanifold{}而言，像 $h_C(F_z)$ 应当等于 $F_z$。类似地，每个 $F_z$ 都需要有一个\trPinStructure{}$P_z^{\#}$，使得 $(h_C|_{F_z})^{*}(P_z^{\#})$ \trIsomorphic{}于 $P_z^{\#}$（该\trIsomorphism{}的一个选取决定了 $F|_C$ 的一个相对\trPinStructure）。
\begin{example}\label{thm:relative-brane-local-model}
取\cref{thm:local-model-sections}中的局部模型，为简单起见取 $r=1$。标准的相对二次复\trVolume{}形式是 $\eta_E=dx_1\wedge\cdots\wedge dx_{n+1}$ 与 $\eta_S=dz$ 的平方商。对每个 $x\in F_z$，
\begin{equation*}
\alpha_{E/S}(TF_x^v)=\frac{\eta_E(v_1,\ldots,v_{n+1})^2}{\eta_S(iz\partial_z)^2},
\end{equation*}
其中 $\{v_1,\ldots,v_{n+1}\}$ 是整个 $F$ 的\trTangentSpace{}的一组\trOrthonormal{}基。具体地，$v_1,\ldots,v_n$ 可以选得与 $\sqrt{z}\,S^n\subset\sqrt{z}\,\RR^{n+1}$ 相切，而 $v_{n+1}=ix\in i\sqrt{z}\,\RR^{n+1}$；于是 $\alpha_{E/S}(TF_x^v)=z^{n-1}$。因此，当 $n=1$ 时 $F$ 容许\trRelative\trGrading，而当 $n>1$ 时则不容许。这是如下事实的一个推论：任何\trDehnTwist{}$\tau_V$，在配备其自然\trGrading{}时，会把中心\trSphere{}$V$ 的\trGrading{}移动 $1-n$ [125, Lemma 5.7]。相比之下，相对\trPinStructure{}总是存在的；事实上，当 $n=1$ 时有四种不同的选取，而当 $n>1$ 时有两种。
\end{example}
\begingroup\tolerance=9999\emergencystretch=3em\hbadness=10000
\trGrading{}的存在保证了如下结论：对任何边界值落在 $F$ 上的\trSection{}$u$，$(u^{*}(TE^{v}),u^{*}(TF^{v}))$ 这一对的相对第一\trChernClass{}为零，从而（在\trRegularity{}假设之下）\trModuliSpace{}的维数将是 $n\chi(S)$。如前那样，在 $\partial S$ 的每个\trConnectedComponent{}上取一点，并把这些点记作 $z_j$。给定一个相对\trPinStructure，我们便得到 $\scr{M}_{E/S}$ 相对于诸 $z_j$ 处\trEvaluationMap{}之积的一个优先\trOrientation。这是\cref{thm:general-surface-index-formula}的一个推论（把抽象\trBraneStructure{}的语言翻译成此处的语言）。带着该优先相对\trOrientation，诸\trEvaluationMap{}于是表示一个\trCohomology{}类（由\trPoincareDuality）
\begin{equation}\label{eq:relative-invariant-class}
\Phi_{E/S,z_1,\ldots,z_k}\in H^{n(k-\chi(S))}\Bigl(\prod_j F_{z_j}\Bigr).
\end{equation}
\endgroup
取某个\trRealLineBundle{}$\beta\to\partial S$。可以这样修改相对\trPinStructure：$P^{\#}\mapsto P^{\#}\otimes(\pi|F)^{*}\beta$；其对 \eqref{eq:relative-invariant-class} 的效果是把该\trInvariant{}乘以 $(-1)^{i}$，其中 $i$ 是 $\beta$ 在其为非平凡的那些\trConnectedComponent{}的数目。最后，这整个讨论可以毫无困难地推广到\trMovingBoundaryCondition{}的情形。
对我们的目的而言，\eqref{eq:relative-invariant-class} 最重要的特例是 $S$ 为\trDisk、而\trFiber{}维数为 $2n=2$ 的情形。后一假设意味着诸\trLinearization{}算子 $D_{E/S,u}$ 都是\trLineBundle{}上的 $\bar\partial$-算子，于是\cref{sub:n-operators-line-bundles}中的论证便可适用。特别地，由 \eqref{eq:winding-number-count} 可见，对每个 $u\in\scr{M}_{E/S}$，$\ker(D_{E/S,u})$ 都是一维的，并且由 $u^{*}(TE^{v})$ 的一个无处为零的\trSection{}$X$ 生成。在由 $X$ 给出的\trTrivialization{}之下，$D_{E/S,u}$ 成为平凡\trLineBundle{}上带常值\trBoundaryCondition{}的标准 $\bar\partial$-算子。回看\cref{thm:abstract-brane-loop-determinant}的证明便可看出：如果 $u^{*}(TF^{v})$ 上的\trPinStructure{}是平凡的，那么相对\trOrientation
$\ker(D_{E/S,u})\cong(TF_z)_{u(z)}$ 便由映射 $X\mapsto X(z)$ 给出；否则，\trOrientation{}取相反的那一个。这意味着，\trLocallyConstant{}函数 $\Phi_{E/S,z}\in H^0(F_z)$ 可以按如下方式计算：取一点 $x\in F_z$。然后对满足 $u(z)=x$ 的诸\trSection{}$u\in\scr{M}_{E/S}$ 作计数：若 $u^{*}(TF^{v})$ 上的\trPinStructure{}平凡，则赋以 $+1$ \trMultiplicity，否则赋以 $-1$。
\begin{example}\label{thm:local-model-n1-evaluation}
我们回到\cref{thm:local-model-sections}，取 $n=1$。该\trBoundaryCondition{}容许一个显式\trParametrization{}$S^1\times S^1\to F$，$(t_1,t_2)\mapsto\bigl(e^{it_1/2}\cos(t_1/2+t_2),\,e^{it_1/2}\sin(t_1/2+t_2)\bigr)$。\trSection{}的\trModuliSpace{}是两个\trCircle{}的\trDisjointUnion{}$\scr{M}_{E/S,\pm}$，它由在 \eqref{eq:local-model-sections} 中取 $a=\tfrac{1}{2}(e^{it},\pm ie^{it})$ 而得到。在这些坐标下，沿边界的\trEvaluationMap{}成为
\begin{equation*}
\begin{aligned}
&\partial S\times\scr{M}_{E/S,+}\longrightarrow F, &&(s,t)\longmapsto(t_1,t_2)=(s,t-s)\\
&\partial S\times\scr{M}_{E/S,-}\longrightarrow F, &&(s,t)\longmapsto(t_1,t_2)=(s,t).
\end{aligned}
\end{equation*}
如果我们在 $F$ 中固定一点，比如 $(t_1,t_2)=(0,0)$，那么每个\trFamily{}中都恰有一条经过该点的\trSection{}$u_{\pm}$，即 $t=0$ 的那一条。注意，在\trHomology{}中，诸\trCircle{}$u_{\pm}(\partial S)\subset F$ 相差一份\trFiber{}$F_z$。鉴于上文所作的一般性观察，这带来如下推论：如果相对\trPinStructure{}在 $F_z$ 上被取成\emph{非平凡}的，那么 $u_+$ 与 $u_-$ 的贡献便相互抵消，从而该\trInvariant{}消没。否则，$\Phi_{E/S,z}=\pm 2$（两种符号都会出现，对应于两个相差 $\partial S$ 上非平凡\trLineBundle{}之\trPullback{}的\trPinStructure）。
\end{example}
取一个带角的\trExactSymplectic\trManifold{}$M$，及其上的一个二次复\trVolume{}形式 $\eta^2_M$。取一个带\trStriplikeEnd{}的\trLefschetzFibration，其在每个端上的\trTrivialization{}都以 $Z^{\pm}\times M$ 为模型。为其配备一个（可以移动的）\trLagrange\trBoundaryCondition{}$F$，并令 $(L_{\xi,0},L_{\xi,1})$ 为与每个端 $\xi$ 相关联的那一对\trLagrange\trSubmanifold。现在设这些\trSubmanifold{}都带有\trBraneStructure{}$(L^{\#}_{\xi,0},L^{\#}_{\xi,1})$，并设 $E$ 带有相对二次复\trVolume{}形式 $\eta^2_{E/S}$，它在每个端上与 $\eta^2_M$ 相符。那么，$F$ 上一个\trCompatible{}的\trRelative\trBraneStructure{}由如下两部分组成：一是相对\trGrading{}$\alpha^{\#}$，它在每个端上等于诸 $L_{\xi,k}$ 的那一个；二是相对\trPinStructure{}$P^{\#}$，连同该结构与诸 $L_{\xi,k}$ 在端上的那些结构之间的一个指定\trIsomorphism。有了这类结构，我们便能对诸\trDeterminantLine{}丛进行与\cref{sub:general-riemann-surfaces}中类似的分析。结果在下述情形中最容易陈述：每条\trBoundaryCircle{}都恰好含一个入向无穷远点（在我们的所有应用中情形都将如此）；此时，对 $\scr{M}_{E/S}(\{y_{\xi}\})$ 中诸\trIsolated{}点的一个带符号计数给出 \eqref{eq:cphi-fibration-map} 的一个更精致的版本，即一个\trCochainComplex{}之间的\trChainMap
\begin{equation*}
C\Phi_{E/S}\colon\bigotimes_{\zeta\in\Sigma^{+}}CF^{*}\bigl(L^{\#}_{\zeta^{+},0},L^{\#}_{\zeta^{+},1}\bigr)\longrightarrow\bigotimes_{\zeta\in\Sigma^{-}}CF^{*}\bigl(L^{\#}_{\zeta^{-},0},L^{\#}_{\zeta^{-},1}\bigr)\bigl[n(\chi(S))-|\Sigma^{-}|\bigr],
\end{equation*}
其中 $CF^{*}$ 是取任意\trCoefficients{}域 $\KK$ 的 $\ZZ$-\trGraded\trFloer\trCochainComplex。更一般的\trSurface{}也是容许的，但那样一来，人们要么需要对某些 $F_z$ 作\trOrientation（见 \eqref{eq:extra-tangent-factors}），要么需要更改与诸端相关联的某些\trBraneStructure（如 \eqref{eq:disc-even-formula} 中那样）。
